% CONCLUSION
% Last generated: 2026-01-15T07:27:20.548Z

\subsection{Summary}

We have presented a comprehensive multi-agent simulation framework for studying DAO governance mechanisms. Our framework supports multiple voting systems, agent heterogeneity, and systematic parameter exploration, enabling reproducible research into fundamental questions of decentralized governance.

Across {{TOTAL_RUNS}} simulation runs and {{EXPERIMENT_COUNT}} experiment sets, we target the following empirical questions:

\begin{enumerate}
    \item \textbf{Quorum sensitivity}: How proposal pass rates change with quorum thresholds under different baseline participation regimes.

    \item \textbf{Scale effects}: Whether participation decays as DAOs scale, and the functional form of that decay.

    \item \textbf{Mechanism trade-offs}: How voting mechanisms redistribute influence and affect decision latency.

    \item \textbf{Design principles}: What simulation-conditional guidance emerges from the above mechanisms and constraints.
\end{enumerate}

\subsection{Contributions}

This work makes the following contributions:

\begin{enumerate}
    \item An open-source simulation framework for DAO governance research
    \item A research CLI enabling reproducible batch experiments
    \item Empirical analysis infrastructure for quorum, scale, and mechanism effects
    \item A theoretical framework connecting multi-agent systems, mechanism design, and social choice
\end{enumerate}

\subsection{Implications}

For researchers, our framework provides infrastructure for systematic governance studies. The combination of formal models, simulation capabilities, and empirical methodology establishes a foundation for computational governance science.

For practitioners, this work provides a structured way to explore parameter space and test governance hypotheses before on-chain deployment. Specific numeric guidance should be interpreted as simulation-conditional and comparative rather than universal prescriptions.

For the broader ecosystem, our work contributes to the maturation of decentralized governance from experimentation to engineering discipline.

\subsection{Closing Remarks}

DAOs represent a significant experiment in organizational design, with potential to reshape how humans coordinate at scale. Understanding their dynamics through rigorous simulation and analysis is essential to realizing this potential while avoiding pitfalls.

We release our framework as open-source software, inviting the community to extend, critique, and improve upon this work. The future of governance is being written now; we hope this contribution helps write it well.

% === LIVING DOCUMENT NOTE ===
% This paper is a living document, updated as new experiments are run.
% Current version reflects {{EXPERIMENT_COUNT}} experiment sets with {{TOTAL_RUNS}} total runs.
% Last updated: {{TIMESTAMP}}
